% Chapter 10: Combinatorial Dimensions
% 3 revelations (ds_dimension, natarajan_dimension, star_number),
% 10 definitions with brief treatment.
% Technical core of Part III. Full Sauer--Shelah proof. Multiclass detective story.

\chapter{Combinatorial Dimensions}
\label{ch:dimensions}

Part~II characterized learnability in four paradigms, and in each case the
answer was a number: the VC dimension for PAC learning, the Littlestone
dimension for online learning, the mind-change ordinal for Gold-style
identification, the teaching dimension for exact learning.  This chapter and
the three that follow develop these numbers (and the two dozen others in the
concept graph) as objects of study in their own right.

The present chapter treats \emph{combinatorial dimensions}: quantities defined
by asking how richly a hypothesis class can label finite configurations of
points.  The simplest such quantity, the VC dimension, asks how many points
can be shattered.  The others refine, extend, or replace this question for
settings beyond binary classification.  We prove the two foundational
combinatorial results in full (the Sauer--Shelah lemma and the
VC dimension characterization already established in \Cref{ch:pac}), recall the
Littlestone dimension from \Cref{ch:online}, and then turn to the chapter's
narrative centerpiece: the thirty-year search for the correct combinatorial
dimension characterizing multiclass learnability, which ended only in 2022.

\medskip

The chapter is organized as follows.

\begin{enumerate}[leftmargin=*,itemsep=2pt]
\item \textbf{VC dimension and shattering} (\Cref{sec:vc-dim}): the
  definitions, the growth function, and the full proof of the Sauer--Shelah
  lemma.
\item \textbf{The Littlestone dimension} (\Cref{sec:ldim-recall}): a brief
  recall from \Cref{ch:online} and the relationship $\VC \leq \Ldim$.
\item \textbf{Beyond binary: pseudodimension and fat-shattering}
  (\Cref{sec:beyond-binary}): scale-sensitive dimensions for real-valued
  function classes.
\item \textbf{The multiclass story: from Natarajan to DS}
  (\Cref{sec:multiclass}): the detective story of multiclass PAC
  characterization.
\item \textbf{Other dimensions} (\Cref{sec:other-dims}): a concise catalog of
  further combinatorial dimensions, with forward references.
\end{enumerate}


% ================================================================
% SECTION 1: VC DIMENSION AND SHATTERING
% ================================================================

\section{VC Dimension and Shattering}
\label{sec:vc-dim}

The definitions in this section were introduced informally in \Cref{ch:objects}
and used without proof in \Cref{ch:pac}.  We now state them precisely and prove
the combinatorial backbone that supports the Fundamental Theorem.

\begin{defn}[Restriction and Shattering]
\label{def:shattering}
Let $\mathcal{H} \subseteq \{0,1\}^X$ be a hypothesis class and let
$S = \{x_1, \ldots, x_n\} \subseteq X$ be a finite set.  The
\emph{restriction} of $\mathcal{H}$ to $S$ is
\[
  \mathcal{H}|_S \;=\; \bigl\{\bigl(h(x_1), \ldots, h(x_n)\bigr) : h \in \mathcal{H}\bigr\}
  \;\subseteq\; \{0,1\}^n.
\]
We say $\mathcal{H}$ \emph{shatters} $S$ if $\mathcal{H}|_S = \{0,1\}^n$,
i.e., every labeling of $S$ is realized by some hypothesis in $\mathcal{H}$.\formal{FLT_Proofs.Complexity.VCDimension.Shatters}
\end{defn}

\begin{defn}[VC Dimension {\cite{VapnikChervonenkis1971}}]
\label{def:vc-dim}
The \emph{Vapnik--Chervonenkis dimension} of $\mathcal{H}$, denoted
$\VC(\mathcal{H})$, is the largest cardinality of a set shattered by
$\mathcal{H}$:
\[
  \VC(\mathcal{H}) \;=\; \max\bigl\{|S| : S \subseteq X,\; \mathcal{H}
  \text{ shatters } S\bigr\}.
\]
If $\mathcal{H}$ shatters arbitrarily large finite sets,
$\VC(\mathcal{H}) = \infty$.\formal{FLT_Proofs.Complexity.VCDimension.VCDim}
\end{defn}

\begin{example}[Halfplanes in $\R^2$]
\label{ex:halfplanes-shatter}
Let $\mathcal{H}$ be the class of halfplanes in $\R^2$:
$\mathcal{H} = \bigl\{x \mapsto \ind[\langle w, x \rangle \geq b] : w \in
\R^2,\, b \in \R\bigr\}$.  Any three non-collinear points are shattered
(\Cref{fig:shatter-halfplane}), but no four points can be: for any four
points, at least one is inside the convex hull of the other three, and the
labeling that assigns it the opposite label from the other three cannot be
realized by a halfplane.  Hence $\VC(\mathcal{H}) = 3$.
\end{example}

\begin{figure}[ht]
\centering
\begin{tikzpicture}[scale=0.7,
    pospt/.style={circle, fill=thmred, inner sep=1.8pt},
    negpt/.style={circle, fill=defblue, inner sep=1.8pt},
    label0/.style={font=\scriptsize, text=defblue},
    label1/.style={font=\scriptsize, text=thmred},
    sep/.style={thick, black!65},
    poshalf/.style={fill=thmred!9},
]
  % Header note: the eight witnesses give VC >= 3
  \node[font=\footnotesize, text=black!78, align=center] at (5.5,2.95)
    {All $2^3=8$ labelings realized by a halfplane: the three points are shattered, so $\VC\ge 3$.};
  % 8 labelings, each WITH its separating halfplane, in 2 rows of 4.
  % fields: col / la / lb / lc / lt(1=horiz,2=vert) / lp(line position) / pd(1=upper|right shaded,0=lower|left)
  \foreach \col/\la/\lb/\lc/\lt/\lp/\pd in {
    0/0/0/0/1/2.0/1,
    1/0/0/1/1/0.85/1,
    2/0/1/0/2/1.5/1,
    3/0/1/1/2/0.5/1,
    4/1/0/0/2/0.5/0,
    5/1/0/1/2/1.5/0,
    6/1/1/0/1/0.85/0,
    7/1/1/1/1/-0.6/1%
  } {
    \pgfmathtruncatemacro{\cx}{mod(\col,4)*3}
    \pgfmathtruncatemacro{\cyr}{int(\col/4)}
    \begin{scope}[shift={(\cx,{-\cyr*3.4})}]
      % separating line + lightly shaded label-1 (positive) halfplane, drawn first (behind points)
      \ifnum\lt=1
        \ifnum\pd=1
          \fill[poshalf] (-0.45,\lp) rectangle (2.45,2.2);
        \else
          \fill[poshalf] (-0.45,-0.9) rectangle (2.45,\lp);
        \fi
        \draw[sep] (-0.45,\lp) -- (2.45,\lp);
      \else
        \ifnum\pd=1
          \fill[poshalf] (\lp,-0.9) rectangle (2.45,2.2);
        \else
          \fill[poshalf] (-0.45,-0.9) rectangle (\lp,2.2);
        \fi
        \draw[sep] (\lp,-0.9) -- (\lp,2.2);
      \fi
      % the three points, on top of the shade
      \ifnum\la=1 \node[pospt] (a\col) at (0,0) {}; \node[label1, below=1pt] at (a\col) {$1$};
      \else        \node[negpt] (a\col) at (0,0) {}; \node[label0, below=1pt] at (a\col) {$0$};\fi
      \ifnum\lb=1 \node[pospt] (b\col) at (2,0) {}; \node[label1, below=1pt] at (b\col) {$1$};
      \else        \node[negpt] (b\col) at (2,0) {}; \node[label0, below=1pt] at (b\col) {$0$};\fi
      \ifnum\lc=1 \node[pospt] (c\col) at (1,1.5) {}; \node[label1, above=1pt] at (c\col) {$1$};
      \else        \node[negpt] (c\col) at (1,1.5) {}; \node[label0, above=1pt] at (c\col) {$0$};\fi
    \end{scope}
  }
  % The d=4 obstruction (Radon): VC < 4, hence VC = 3.
  \begin{scope}[shift={(4.0,-7.9)}]
    \coordinate (p1) at (0,0);
    \coordinate (p2) at (3,0);
    \coordinate (p3) at (1.5,2.2);
    \coordinate (pin) at (1.5,0.75);
    \fill[black!6] (p1) -- (p2) -- (p3) -- cycle;
    \draw[black!40, dashed] (p1) -- (p2) -- (p3) -- cycle;
    \node[negpt] at (p1) {}; \node[label0, below=1pt] at (p1) {$0$};
    \node[negpt] at (p2) {}; \node[label0, below=1pt] at (p2) {$0$};
    \node[negpt] at (p3) {}; \node[label0, above=1pt] at (p3) {$0$};
    \node[pospt] at (pin) {}; \node[label1, right=1pt] at (pin) {$1$};
    \node[font=\footnotesize, text=black!78, align=center] at (1.5,-1.2)
      {No four points are shattered: the inner point lies in the hull of the others, so the\\labeling isolating it is unrealizable by a halfplane (Radon). Hence $\VC<4$, so $\VC=3$.};
  \end{scope}
\end{tikzpicture}
\caption[Shattering three points: each of the eight labelings realized by a separating halfplane, and no four points shattered, so the VC dimension is three.]{The eight labelings of three non-collinear points, each drawn with a separating halfplane that realizes it: the boundary splits the label-one points (red) from the label-zero points (blue), and the label-one side is lightly shaded.  A witness line in every cell means the restriction of the halfplane class to these three points is the full set of eight labelings, so the points are shattered (\leanname{Shatters}) and the VC dimension is at least three.  The lower panel shows why no four points are shattered: one of any four lies in the convex hull of the others, and the labeling isolating it is unrealizable by a halfplane (Radon), so the VC dimension (\leanname{VCDim}) equals three, which is $d+1$ for halfspaces in $d$ dimensions.  The shattering and VC-dimension concepts are the verified definitions; the value three and the four-point obstruction are classical.}
\label{fig:shatter-halfplane}
\end{figure}


\subsection{The Growth Function}

The growth function measures how many distinct labelings $\mathcal{H}$ can
produce on sets of a given size.

\begin{defn}[Growth Function]
\label{def:growth-function}
The \emph{growth function} (or \emph{shatter coefficient}) of $\mathcal{H}$ is
\[
  \growth_\mathcal{H}(n) \;=\; \max_{S \subseteq X,\; |S|=n}
  \bigl|\mathcal{H}|_S\bigr|.
\]
Equivalently, $\growth_\mathcal{H}(n)$ is the maximum number of distinct
behaviors $\mathcal{H}$ can exhibit on any $n$ points.\formal{FLT_Proofs.Complexity.VCDimension.GrowthFunction}
\end{defn}

The growth function satisfies $\growth_\mathcal{H}(n) \leq 2^n$ always, with
equality when $\mathcal{H}$ shatters some set of size~$n$.  The remarkable
content of the Sauer--Shelah lemma is that once shattering fails (at
$n = \VC(\mathcal{H}) + 1$), the growth function drops from exponential to
polynomial \emph{immediately and permanently}.


\subsection{The Sauer--Shelah Lemma}

\begin{lemma}[Sauer--Shelah {\cite{Sauer1972,Shelah1972}}]
\label{lem:sauer-shelah}
Let $\mathcal{H} \subseteq \{0,1\}^X$ with $\VC(\mathcal{H}) = d$.  Then for
all $n \geq 1$,
\[
  \growth_\mathcal{H}(n) \;\leq\; \sum_{i=0}^{d} \binom{n}{i}.
\]
In particular, $\growth_\mathcal{H}(n) \leq \bigl(\frac{en}{d}\bigr)^d$ for
$n \geq d \geq 1$, so the growth function is polynomial in $n$ when
$d < \infty$.\formal{FLT_Proofs.Theorem.PAC.sauer_shelah_quantitative}
\end{lemma}

The proof is by double induction and is one of the most elegant arguments in
combinatorics.  The key idea is to split the hypothesis class along one
coordinate and apply the inductive hypothesis to the two resulting classes,
whose VC dimensions are carefully controlled.

\begin{proof}
\label{lem:sauer-shelah-proof}
We prove a stronger statement: for any finite set $S$ with $|S| = n$ and
any class $\mathcal{H}' \subseteq \{0,1\}^S$ (not necessarily arising from
a restriction),
\[
  |\mathcal{H}'| \;\leq\; \bigl|\{A \subseteq S : \mathcal{H}' \text{ shatters } A\}\bigr|.
  \tag{$\star$}
\]
This implies the lemma: if $\VC(\mathcal{H}') \leq d$, then the right-hand
side counts only subsets $A$ with $|A| \leq d$, of which there are at most
$\sum_{i=0}^{d}\binom{n}{i}$.

We prove $(\star)$ by induction on $|S|$.

\smallskip\noindent
\textbf{Base case} ($|S| = 1$, say $S = \{x\}$).  If
$\mathcal{H}' = \{0\}$ or $\mathcal{H}' = \{1\}$, then $|\mathcal{H}'| = 1$
and $\mathcal{H}'$ shatters $\emptyset$, giving at least one shattered set.
If $\mathcal{H}' = \{0,1\}$, then $|\mathcal{H}'| = 2$ and $\mathcal{H}'$
shatters both $\emptyset$ and $\{x\}$, giving at least two shattered sets.
In both cases $(\star)$ holds.

\smallskip\noindent
\textbf{Inductive step.}  Let $|S| = n > 1$.  Pick any element $x \in S$
and set $S' = S \setminus \{x\}$.  Partition $\mathcal{H}'$ according to
behavior on $x$:
\begin{align*}
  \mathcal{H}_0 &= \{h|_{S'} : h \in \mathcal{H}',\; h(x) = 0\}, \\
  \mathcal{H}_1 &= \{h|_{S'} : h \in \mathcal{H}',\; h(x) = 1\}, \\
  \mathcal{H}_\cap &= \mathcal{H}_0 \cap \mathcal{H}_1.
\end{align*}
Here $\mathcal{H}_\cap$ consists of the functions on $S'$ that appear with
\emph{both} labels on $x$.

Count: each $h \in \mathcal{H}'$ contributes to exactly one of $\mathcal{H}_0$
or $\mathcal{H}_1$ (as a function on $S'$), but a function in
$\mathcal{H}_\cap$ is contributed by two elements of $\mathcal{H}'$ (one
with $h(x)=0$, one with $h(x)=1$).  Therefore,
\[
  |\mathcal{H}'| \;=\; |\mathcal{H}_0| + |\mathcal{H}_1| - |\mathcal{H}_\cap|
  + |\mathcal{H}_\cap|
  \;=\; |\mathcal{H}_0 \cup \mathcal{H}_1| + |\mathcal{H}_\cap|.
\]
More precisely, $|\mathcal{H}'| = |\mathcal{H}_0| + |\mathcal{H}_1|$, since
functions in $\mathcal{H}_0 \setminus \mathcal{H}_\cap$ appear only from
$h(x) = 0$, functions in $\mathcal{H}_1 \setminus \mathcal{H}_\cap$ appear
only from $h(x) = 1$, and functions in $\mathcal{H}_\cap$ appear twice.  So
\[
  |\mathcal{H}'| = |\mathcal{H}_0 \cup \mathcal{H}_1| + |\mathcal{H}_\cap|.
\]

Apply the inductive hypothesis (on $S'$, which has $n-1$ elements) to both
$\mathcal{H}_0 \cup \mathcal{H}_1$ and $\mathcal{H}_\cap$:
\begin{align*}
  |\mathcal{H}_0 \cup \mathcal{H}_1|
    &\leq \bigl|\{A \subseteq S' : (\mathcal{H}_0 \cup \mathcal{H}_1) \text{ shatters } A\}\bigr|, \\
  |\mathcal{H}_\cap|
    &\leq \bigl|\{A \subseteq S' : \mathcal{H}_\cap \text{ shatters } A\}\bigr|.
\end{align*}

Now we connect the shattered sets of the restricted classes back to those
of $\mathcal{H}'$:

\begin{enumerate}[leftmargin=*,itemsep=3pt]
\item If $A \subseteq S'$ is shattered by $\mathcal{H}_0 \cup \mathcal{H}_1$
  but \emph{not} by $\mathcal{H}_\cap$, then $A$ is shattered by $\mathcal{H}'$
  as well (since every labeling of $A$ is realized by some $h|_{S'} \in
  \mathcal{H}_0 \cup \mathcal{H}_1$, and at least one such $h$ exists in
  $\mathcal{H}'$).

\item If $A \subseteq S'$ is shattered by $\mathcal{H}_\cap$, then every
  labeling of $A$ is realized by functions appearing with both labels on $x$.
  Hence $\mathcal{H}'$ shatters $A \cup \{x\}$: for any labeling of
  $A \cup \{x\}$, pick the label $b$ assigned to $x$, then pick $h|_{S'} \in
  \mathcal{H}_\cap$ realizing the labeling of $A$; the corresponding
  $h \in \mathcal{H}'$ with $h(x) = b$ realizes the full labeling.
\end{enumerate}

The sets in (1) do not contain $x$, and the sets in (2) do
contain $x$, so they are disjoint families.  Therefore,
\[
  |\mathcal{H}'|
  \;\leq\; \bigl|\{A \subseteq S : \mathcal{H}' \text{ shatters } A,\;
  x \notin A\}\bigr|
  + \bigl|\{A \subseteq S : \mathcal{H}' \text{ shatters } A,\;
  x \in A\}\bigr|
  \;=\; \bigl|\{A \subseteq S : \mathcal{H}' \text{ shatters } A\}\bigr|.
\]
This completes the induction.
\end{proof}

\begin{remark}
The lemma was proved independently by Sauer~\cite{Sauer1972} and
Shelah~\cite{Shelah1972}, with a related result by Perles and Shelah
appearing earlier.  Vapnik and Chervonenkis~\cite{VapnikChervonenkis1971}
proved a weaker bound in their 1971 paper.  The proof above follows the
elegant ``number of shattered sets'' formulation due to Alon and
Frankl.
\end{remark}

\begin{figure}[ht]
\centering
\begin{tikzpicture}
\begin{axis}[
    width=12cm, height=7cm,
    xlabel={$n$ (number of points)},
    ylabel={Number of labelings $\growth_\mathcal{H}(n)$},
    ymode=log, log basis y=10,
    xmin=0, xmax=16, ymin=1, ymax=1e5,
    xtick={0,2,4,6,8,10,12,14,16},
    ytick={1,10,100,1000,10000,100000},
    legend pos=south east,
    legend cell align={left},
    legend style={font=\small, fill=white, fill opacity=0.9, draw=gray!40, text opacity=1},
    grid=major,
    grid style={gray!20},
    thick,
    clip mode=individual,
]
% 2^n reference ray: the ceiling, attained while the class still shatters (dashed)
\addplot[domain=0:16, samples=120, thick, thmred, dashed] {2^x};
\addlegendentry{$2^n$ (shatters forever)}

% Actual growth of a maximum VC=3 class = the Sauer--Shelah bound:
% equals 2^n for n<=3 (binomials i>n vanish), then peels permanently below (solid, integer marks)
\addplot[domain=0:16,
         samples at={0,1,2,3,4,5,6,7,8,9,10,11,12,13,14,15,16},
         thick, defblue, mark=*, mark size=1pt]
    {1 + x + x*(x-1)/2 + x*(x-1)*(x-2)/6};
\addlegendentry{$\growth_\mathcal{H}(n)=\sum_{i=0}^{3}\binom{n}{i}$ \;($\VC=3$)}

% Marker: n = d = 3 (last shattered size)
\draw[thick, dashed, notegray] (axis cs:3,1) -- (axis cs:3,1e5);
\node[font=\scriptsize, text=notegray, anchor=south, rotate=90]
    at (axis cs:3,5e2) {$n=d=3$ (last shattered)};

% Marker: n = d+1 = 4 (Sauer--Shelah bound bites; permanent drop)
\draw[thick, dashed, notegray] (axis cs:4,1) -- (axis cs:4,1e5);
\node[font=\scriptsize, text=notegray, anchor=south, rotate=90]
    at (axis cs:4,5e2) {$n=d{+}1$: bound bites};

% The widening exponential gap for n > d
\draw[<->, thin, gray!70] (axis cs:11,2048) -- (axis cs:11,232);
\node[font=\scriptsize, text=gray!75, anchor=west]
    at (axis cs:11.3,720) {exponential saving};

\end{axis}
\end{tikzpicture}
\caption[The growth function: exponential while the class shatters, then polynomial of degree~$\VC$.]{The
growth function $\growth_\mathcal{H}(n)$ (\leanname{GrowthFunction}) counts the most
distinct labelings a class can produce on any $n$ points.  It never exceeds $2^n$ (the
dashed ray) and equals $2^n$ exactly while the class still shatters a set of that size,
that is for $n \leq d = \VC(\mathcal{H})$.  The solid curve is the largest a class of VC
dimension three can grow, the Sauer--Shelah bound $\sum_{i=0}^{3}\binom{n}{i}$
(\leanname{sauer_shelah_quantitative}): it rides the $2^n$ ray through $n=3$, then at
$n = d+1 = 4$ peels permanently away into a polynomial arc whose degree is the VC
dimension.  The drop is immediate and there is no intermediate growth rate.  On the
logarithmic vertical axis the exponential $2^n$ is a straight line and the polynomial
bends below it, so the widening gap is the exponential saving a finite VC dimension buys.
This collapse is the combinatorial step the Fundamental Theorem turns into PAC
learnability, since the growth function is the quantity a uniform-convergence bound
controls.  The definition and the bound are both verified; the asymptotic form
$(en/d)^d$ and the no-middle-regime reading are classical.}
\label{fig:growth-function}
\end{figure}

The Sauer--Shelah lemma is the combinatorial engine behind the Fundamental
Theorem of Statistical Learning (\Cref{thm:fundamental}): it converts
a finite VC dimension into a polynomial bound on the growth function, which in
turn yields uniform convergence of empirical risks via a symmetrization
argument.


% ================================================================
% SECTION 2: THE LITTLESTONE DIMENSION (CROSS-REFERENCE)
% ================================================================

\section{The Littlestone Dimension}
\label{sec:ldim-recall}

The VC dimension measures shatterability of \emph{sets}: it asks for the
largest set whose every labeling is realized.  The Littlestone dimension,
introduced in \Cref{ch:online} and fully characterized there, measures
shatterability of \emph{trees}: it asks for the deepest complete binary tree
whose every root-to-leaf path is realized.  We recall the definition for
reference.

\begin{defn}[Mistake Tree, Recall]
\label{def:mistake-tree-recall}
A \emph{mistake tree} for $\mathcal{H} \subseteq \{0,1\}^X$ is a complete
binary tree in which each internal node is labeled with an instance $x \in X$,
left edges correspond to label~$0$, right edges to label~$1$, and every
root-to-leaf path is consistent with some $h \in \mathcal{H}$
(see \Cref{def:mistake-tree}).
\end{defn}

\begin{defn}[Littlestone Dimension, Recall]
\label{def:ldim-recall}
The \emph{Littlestone dimension} $\Ldim(\mathcal{H})$ is the maximum depth
of a mistake tree for $\mathcal{H}$ (\Cref{def:ldim}).  The
Littlestone characterization (\Cref{thm:littlestone-characterization}) establishes that
$\Ldim(\mathcal{H})$ equals the optimal worst-case mistake bound.
\end{defn}

The relationship between the two fundamental dimensions is strict:

\begin{prop}
\label{prop:vc-leq-ldim-dim}
For any hypothesis class $\mathcal{H}$,
$\VC(\mathcal{H}) \leq \Ldim(\mathcal{H})$.\formal{FLT_Proofs.Complexity.Littlestone.BranchWiseLittlestoneDim_ge_VCDim}
\end{prop}

\begin{proof}
If $\mathcal{H}$ shatters $\{x_1, \ldots, x_d\}$, one can construct a
mistake tree of depth $d$ by placing $x_i$ at every node of depth $i$:
since every labeling is realized, every root-to-leaf path is consistent with
some $h \in \mathcal{H}$.  Hence any shattered set yields a mistake tree of
the same depth.
\end{proof}

The converse fails: threshold classifiers on $\R$ have $\VC = 1$ but
$\Ldim = \infty$ (\Cref{ch:separations}).  Intuitively, the Littlestone
dimension is larger because it accounts for adaptive adversaries: the
adversary in the online game chooses instances sequentially, seeing the
learner's responses, while the VC dimension considers all labelings of a
fixed set simultaneously.


% ================================================================
% SECTION 3: BEYOND BINARY
% ================================================================

\section{Beyond Binary: Pseudodimension and Fat-Shattering}
\label{sec:beyond-binary}

When the hypothesis class consists of real-valued functions
$\mathcal{F} \subseteq \R^X$ rather than binary classifiers, shattering must
be redefined.  Two progressively refined notions capture the relevant
complexity.

\begin{defn}[Pseudodimension]
\label{def:pseudodim}
Let $\mathcal{F} \subseteq \R^X$.  A set $S = \{x_1, \ldots, x_n\} \subseteq X$
is \emph{pseudo-shattered} (or \emph{P-shattered}) by $\mathcal{F}$ if there
exist thresholds $t_1, \ldots, t_n \in \R$ such that, for every labeling
$b \in \{0,1\}^n$, there exists $f \in \mathcal{F}$ with
\[
  \mathrm{sgn}\bigl(f(x_i) - t_i\bigr) = (-1)^{1 - b_i}
  \quad \text{for all } i = 1, \ldots, n.
\]
The \emph{pseudodimension} $\Pdim(\mathcal{F})$ is the largest cardinality of
a P-shattered set.\formal{FLT_Proofs.Complexity.Structures.Pseudodimension}
\end{defn}

The thresholds $t_i$ act as ``witnesses'': they reduce a real-valued
shattering problem to a binary one.  When $\mathcal{F}$ is already
$\{0,1\}$-valued, the thresholds $t_i = 1/2$ recover the standard VC
dimension, so $\Pdim$ is a proper generalization.

\begin{defn}[Fat-Shattering Dimension {\cite{AlonBenDavidCesaBianchiHaussler1997}}]
\label{def:fat-shattering}
Let $\mathcal{F} \subseteq \R^X$ and $\gamma > 0$.  A set
$S = \{x_1, \ldots, x_n\}$ is \emph{$\gamma$-fat-shattered} by $\mathcal{F}$
if there exist thresholds $t_1, \ldots, t_n \in \R$ such that, for every
labeling $b \in \{0,1\}^n$, there exists $f \in \mathcal{F}$ with
\[
  f(x_i) \;\geq\; t_i + \gamma \;\;\text{when } b_i = 1,
  \qquad
  f(x_i) \;\leq\; t_i - \gamma \;\;\text{when } b_i = 0,
\]
for all $i$.  The \emph{fat-shattering dimension at scale $\gamma$}, denoted
$\fatshat_\gamma(\mathcal{F})$, is the largest cardinality of a
$\gamma$-fat-shattered set.\formal{FLT_Proofs.Complexity.Structures.FatShatteringDim}
\end{defn}

The margin parameter $\gamma$ makes fat-shattering \emph{scale-sensitive}:
large margins are harder to achieve, so
$\fatshat_\gamma(\mathcal{F}) \leq \fatshat_{\gamma'}(\mathcal{F})$ whenever
$\gamma \geq \gamma'$.  As $\gamma \to 0$, the fat-shattering dimension
approaches the pseudodimension.

\begin{thm}[Characterization of Real-Valued Learnability
  {\cite{AlonBenDavidCesaBianchiHaussler1997}}]
\label{thm:fat-char}
A real-valued function class $\mathcal{F} \subseteq [0,1]^X$ is
\emph{agnostically PAC learnable} (with respect to the absolute loss) if and
only if $\fatshat_\gamma(\mathcal{F}) < \infty$ for every $\gamma > 0$.
\end{thm}

This is the real-valued analogue of the Fundamental Theorem: just as finite VC
dimension characterizes binary PAC learnability, finite fat-shattering
dimension at every scale characterizes real-valued learnability.  The
pseudodimension alone is not sufficient for characterization in the agnostic
setting, as it ignores the scale structure.

\begin{remark}[Dimension hierarchy]
For a class $\mathcal{F} \subseteq [0,1]^X$:
\[
  \fatshat_\gamma(\mathcal{F})
  \;\leq\; \Pdim(\mathcal{F})
  \;\leq\; \VC\!\bigl(\mathrm{subgraph}(\mathcal{F})\bigr)
\]
where $\mathrm{subgraph}(\mathcal{F}) = \{(x,t) \mapsto \ind[f(x) \geq t] :
f \in \mathcal{F}\}$.  When $\mathcal{F}$ is $\{0,1\}$-valued, all three
quantities coincide with $\VC(\mathcal{F})$.
\end{remark}


% ================================================================
% SECTION 4: THE MULTICLASS STORY
% ================================================================

\section{The Multiclass Story: From Natarajan to DS}
\label{sec:multiclass}

Binary classification has a clean answer: a class $\mathcal{H} \subseteq
\{0,1\}^X$ is PAC learnable if and only if $\VC(\mathcal{H}) < \infty$.
What is the corresponding answer when the label set $Y$ has $k \geq 3$
elements?  This question, first posed in the 1980s, was open for thirty
years.  Its resolution required not merely a new proof technique but a
fundamentally new combinatorial dimension, because the obvious candidate turned out to be wrong.


\subsection{The Natarajan Dimension: The Obvious Candidate}

\begin{defn}[Natarajan Dimension]
\label{def:natarajan}
Let $\mathcal{H} \subseteq Y^X$ with $|Y| \geq 2$.  The \emph{Natarajan
dimension} $\Ndim(\mathcal{H})$ is the largest $d$ such that there exist
\begin{itemize}[leftmargin=*,itemsep=1pt]
\item a set $S = \{x_1, \ldots, x_d\} \subseteq X$, and
\item two functions $f_0, f_1 \colon S \to Y$ with $f_0(x_i) \neq f_1(x_i)$
  for all $i$,
\end{itemize}
such that for every labeling $b \in \{0,1\}^d$, there exists $h \in
\mathcal{H}$ with $h(x_i) = f_{b_i}(x_i)$ for all $i$.\formal{FLT_Proofs.Complexity.Structures.NatarajanDim}
\end{defn}

The Natarajan dimension generalizes the VC dimension in a natural way: instead
of requiring all $2^d$ binary labelings, it requires all $2^d$ labelings
drawn from two specified ``opposing'' functions.  When $|Y| = 2$, the
Natarajan dimension equals the VC dimension exactly.

In the 1980s and 1990s, the working conjecture was:
\begin{quote}
\itshape A multiclass hypothesis class $\mathcal{H} \subseteq Y^X$ is PAC
learnable if and only if $\Ndim(\mathcal{H}) < \infty$.
\end{quote}
This conjecture was supported by partial results.  Ben-David, Cesa-Bianchi,
Haussler, and Long proved that finite Natarajan dimension is
\emph{sufficient} for multiclass PAC learnability when combined with a bound
on $|Y|$.  The ``if'' direction seemed solid.  The conjecture survived for
decades without a counterexample.

It was wrong.


\subsection{The Counterexample: When the Obvious Fails}

\begin{thm}[Failure of Natarajan Dimension
  {\cite{DanielyShalevShwartz2014}}]
\label{thm:natarajan-failure}
There exists a hypothesis class $\mathcal{H} \subseteq Y^X$ with
$\Ndim(\mathcal{H}) = 1$ that is \emph{not} PAC learnable.
\end{thm}

This is a dramatic failure: the Natarajan dimension can be as small as
possible (just one) and still the class can be unlearnable.  The problem lies
in the label set: the construction uses $|Y| = \infty$, and the Natarajan
dimension, which examines only pairs of opposing functions, fails to capture
the combinatorial explosion that arises when the label set is large.

The witness is a class constructed via a combinatorial argument.  The key
insight is that the Natarajan dimension constrains the number of ``binary
slices'' through the labeling space but does not constrain the total
complexity when the label set grows without bound.

\begin{separation}
\textbf{Natarajan vs.\ PAC learnability.}  The Natarajan dimension does not
characterize multiclass PAC learnability:
\[
  \Ndim(\mathcal{H}) < \infty
  \quad\not\!\!\!\implies\quad
  \mathcal{H} \text{ is PAC learnable}.
\]
The gap is witnessed by a class with $\Ndim = 1$, $|Y| = \infty$, that
requires unbounded sample complexity.  The correct characterization requires
the DS dimension (\Cref{def:ds-dim} below).
\end{separation}


\subsection{The DS Dimension: The Correct Answer}

\begin{defn}[DS Dimension {\cite{BrukhimCarmonDinurMoranYehudayoff2022}}]
\label{def:ds-dim}
Let $\mathcal{H} \subseteq Y^X$.  The \emph{DS dimension} (named for
Daniely and Shalev-Shwartz) $\DSdim(\mathcal{H})$ is defined via
pseudo-cubes in the one-inclusion hypergraph of $\mathcal{H}$.

Construct the \emph{one-inclusion hypergraph} $G_\mathcal{H}$ on a finite set
$S = \{x_1, \ldots, x_n\} \subseteq X$ as follows: the vertices are the
hypotheses $\mathcal{H}|_S$ (restrictions of $\mathcal{H}$ to $S$), and for
each $i \in \{1, \ldots, n\}$, add a hyperedge connecting all hypotheses
that agree on $S \setminus \{x_i\}$ (i.e., they differ only at $x_i$).

A \emph{pseudo-cube of dimension $d$} in $G_\mathcal{H}$ is a subhypergraph
isomorphic to the complete $d$-dimensional hypercube graph $\{0,1\}^d$.

The DS dimension $\DSdim(\mathcal{H})$ is the supremum over all finite
$S \subseteq X$ of the maximum dimension of a pseudo-cube in $G_\mathcal{H}$.\footnote{The genuine DS dimension is not machine-checked. The kernel types a related but distinct notion: the \emph{all-labelings} (strong shattering) dimension \leanname{FLT_Proofs.Complexity.Structures.StrongShatteringDim}, the largest $S$ on which every $f \colon S \to Y$ is realizable, whose value is at most $\Ndim$, the reverse of the genuine DS relationship $\Ndim \le \DSdim$. Formalizing the pseudo-cube DS dimension is \cref{op:ds-formal}.}
\end{defn}

\begin{thm}[Multiclass Characterization
  {\cite{BrukhimCarmonDinurMoranYehudayoff2022}}]
\label{thm:ds-char}
A hypothesis class $\mathcal{H} \subseteq Y^X$ is PAC learnable if and only
if $\DSdim(\mathcal{H}) < \infty$.
\end{thm}

This theorem, proved by Brukhim, Chepoi, Daniel, Moran, and Yehudayoff in
2022, resolved the multiclass characterization problem.  It is the
multiclass analogue of the Fundamental Theorem: the DS dimension plays
exactly the role that the VC dimension plays for binary classification.

\begin{historical}
\textbf{Why the Natarajan dimension fails and the DS dimension succeeds.}
The Natarajan dimension examines the labeling structure of $\mathcal{H}$ by
projecting onto pairs of opposing functions, a fundamentally binary lens
applied to a multiclass problem.  The DS dimension, by contrast, examines
the full combinatorial structure of the one-inclusion hypergraph, which
encodes all the ways hypotheses can disagree on single points.

The one-inclusion hypergraph is not new: Haussler, Littlestone, and Warmuth
introduced it in 1994 for binary classification, where it yields an optimal
PAC learner (the one-inclusion algorithm).  What Brukhim et al.\ discovered
is that the right way to measure the complexity of a multiclass hypothesis
class is to look for high-dimensional pseudo-cubes inside this hypergraph.
Pseudo-cubes are the multiclass analogue of shattered sets: they represent
configurations where the labeling structure is rich enough to prevent
learning.
\end{historical}

The dimension hierarchy for multiclass problems is:

\begin{figure}[ht]
\centering
\begin{tikzpicture}[
    dimbox/.style={draw, thick, rounded corners=3pt, minimum width=2.6cm,
      minimum height=0.95cm, font=\small\bfseries, align=center},
    dimboxdash/.style={dimbox, dashed},
    ladsolid/.style={-{Stealth[length=2.5mm]}, thick},
    laddash/.style={-{Stealth[length=2.5mm]}, thick, dashed, draw=notegray},
    branch/.style={thick, defblue},
    solidcall/.style={thick, draw=defblue},
    dashcall/.style={thick, dashed, draw=notegray},
    failcall/.style={thick, dashed, draw=thmred},
    implabel/.style={font=\scriptsize, text=notegray, align=center},
    gaplabel/.style={font=\scriptsize\itshape, text=thmred, align=center},
]
  % --- the magnitude ladder spine: verified floor -> literature ceiling ---
  \node[dimbox, fill=teal!14] (sst) at (0,0) {Strong\\shattering\\{\scriptsize dim.}};
  \node[dimbox, fill=edgeorange!14] (nat) at (4,0) {$\Ndim$\\{\scriptsize Natarajan}};
  \node[dimboxdash, fill=notegray!8] (ds) at (8,0) {$\DSdim$\\{\scriptsize DS / pseudo-cube}};
  % --- VC: the binary face of Ndim, above the spine ---
  \node[dimbox, fill=defblue!14, minimum width=2.2cm] (vc) at (4,2.0) {$\VC$\\{\scriptsize binary}};

  % --- ladder edges: solid = kernel-verified, dashed = literature ---
  \draw[ladsolid] (sst) -- node[implabel, above] {$\le$} (nat);
  \draw[laddash] (nat) -- node[implabel, above] {$\le$} (ds);
  \draw[branch] (vc) -- node[implabel, right, text=defblue] {\scriptsize $|Y|{=}2$:\ $\Ndim{=}\VC$} (nat);

  % --- characterization call-outs (each verdict hangs off its node) ---
  \node[implabel, text=defblue, above=0.55cm of vc] (vcnote)
    {$\VC<\infty \Leftrightarrow$ learnable\\{\scriptsize binary, verified}};
  \draw[solidcall] (vc) -- (vcnote);

  \node[gaplabel, below=0.6cm of nat] (gap)
    {$\Ndim<\infty \not\Rightarrow$ learnable\\{\scriptsize witness $\Ndim=1$, $|Y|=\infty$}\\{\scriptsize (Daniely--Shalev-Shwartz 2014)}};
  \draw[failcall] (nat) -- (gap);

  \node[implabel, below=0.6cm of ds] (dscorrect)
    {$\DSdim<\infty \Leftrightarrow$ learnable\\{\scriptsize (Brukhim et al.\ 2022,}\\{\scriptsize not formalized)}};
  \draw[dashcall] (ds) -- (dscorrect);

  % --- line-style legend: the evidence-grade key ---
  \begin{scope}[shift={(-0.3,-3.75)}]
    \draw[ladsolid] (0,0) -- (0.7,0);
    \node[font=\scriptsize, anchor=west] at (0.85,0) {kernel-verified (sorry-free)};
    \draw[laddash] (5.1,0) -- (5.8,0);
    \node[font=\scriptsize, anchor=west] at (5.95,0) {literature / not yet formalized};
  \end{scope}
\end{tikzpicture}
\caption[The multiclass dimension ladder: verified rungs versus literature.]{The
multiclass dimension ladder, drawn at its evidence grade.  For binary classes
($|Y|=2$) the VC dimension characterizes PAC learnability completely: a class is
learnable if and only if its VC dimension is finite, and this equivalence is
machine-checked (solid).  The obvious generalization to $k \ge 3$ labels is the
Natarajan dimension (\leanname{NatarajanDim}), which equals the VC dimension when
$|Y|=2$; but finite Natarajan dimension does \emph{not} imply learnability
(Daniely--Shalev-Shwartz 2014, with a witness of Natarajan dimension one over an
infinite label set), so the obvious dimension fails to characterize the multiclass
regime.  The genuine characterization is the DS dimension, defined through
pseudo-cubes in the one-inclusion hypergraph (Brukhim et al.\ 2022); it is correct
but not yet formalized, so it and its equivalence are drawn dashed.  In its place the
kernel verifies the all-labelings strong-shattering dimension
(\leanname{StrongShatteringDim}) and the inequality
\leanname{strongShatteringDim_le_NatarajanDim}, which is the \emph{reverse} of the
genuine relationship $\Ndim \le \DSdim$.  Hence the full ladder
$\mathrm{StrongShattering} \le \Ndim \le \DSdim$ has only its lower rung
machine-checked: solid edges are kernel-verified, dashed edges are literature or not
yet formalized.}
\label{fig:dimension-hierarchy}
\end{figure}


% ================================================================
% SECTION 5: OTHER DIMENSIONS
% ================================================================

\section{Other Dimensions}
\label{sec:other-dims}

Several further combinatorial dimensions appear in the concept graph.  We
collect them here in catalog form; each is developed more fully in the chapter
where it plays a characterizing role.

\begin{table}[ht]
\centering
\small
\caption{Catalog of combinatorial dimensions beyond VC, Littlestone,
pseudodimension, fat-shattering, Natarajan, and DS.}
\label{tab:other-dims}
\begin{tabularx}{\textwidth}{@{}l l X l@{}}
\toprule
\textbf{Dimension} & \textbf{Notation} & \textbf{What it measures} & \textbf{Reference} \\
\midrule
Star number
  & $\mathfrak{s}(\mathcal{H})$
  & The largest set on which each point can be isolated by some concept.
    It satisfies $\mathfrak{s}(\mathcal{H}) \geq \VC(\mathcal{H})$, so finite
    star number implies finite VC, but \emph{not} conversely (singletons:
    $\mathfrak{s} = \infty$, $\VC = 1$).  It governs active-learning label
    complexity.
  & \Cref{ch:objects} \\
\addlinespace
Eluder dimension
  & $\dim_E(\mathcal{F})$
  & The length of the longest sequence of points such that each point is
    ``independent'' of the previous ones, given the function class.
    Measures exploration complexity in sequential decision problems.
  & \Cref{ch:extensions} \\
\addlinespace
SQ dimension
  & $\SQdim(\mathcal{C})$
  & The largest set of nearly uncorrelated concepts.  Determines the number
    of statistical queries needed to learn $\mathcal{C}$ in the statistical
    query model, where the learner receives expectations
    $\E[\phi(x,y)]$ rather than individual examples.
  & \Cref{ch:extensions} \\
\addlinespace
VCL tree
  & (none)
  & A tree structure used in the trichotomy theorem (\Cref{ch:universal})
    to separate exponential from arbitrarily slow learners.  Combines VC
    and Littlestone structure.
  & \Cref{ch:universal} \\
\addlinespace
Ordinal VC dim
  & $\VC^\alpha$
  & A transfinite refinement of VC dimension, using ordinal-indexed
    Cantor--Bendixson ranks on the space of consistent hypotheses.
  & \Cref{ch:ordinals} \\
\addlinespace
Ordinal Littlestone dim
  & $\Ldim^\alpha$
  & A transfinite refinement of Littlestone dimension, paralleling
    ordinal VC dimension for the online setting.
  & \Cref{ch:ordinals} \\
\bottomrule
\end{tabularx}
\end{table}

\begin{remark}[Star number versus VC dimension]
\label{rmk:star-number}
It is tempting to believe that finite star number and finite VC dimension
coincide.  They do not.  The star number dominates the VC dimension,
$\mathfrak{s}(\mathcal{H}) \geq \VC(\mathcal{H})$ (\Cref{prop:star-ge-vc}),
because a shattered set is in particular a set on which each point can be
isolated; but the reverse fails, and fails infinitely.  The class of
singletons on an infinite domain has $\VC = 1$ yet $\mathfrak{s} = \infty$
(\Cref{rmk:star-not-equiv}).\formal{FLT_Proofs.Complexity.VCDimension.StarNumber} Finiteness of the star number is therefore
\emph{strictly stronger} than PAC learnability, not equivalent to it; its
proper role is in active learning, where it controls the label
complexity~\cite{HannekeYang2015}.
\end{remark}


% ================================================================
% CLOSING REMARKS
% ================================================================

\bigskip
\begin{center}
\rule{0.5\textwidth}{0.4pt}
\end{center}
\bigskip

The combinatorial dimensions of this chapter are the coordinates of a map.
The VC dimension locates a class in the PAC landscape; the Littlestone
dimension locates it in the online landscape; the DS dimension locates it in
the multiclass landscape; the fat-shattering dimension locates it in the
real-valued landscape.  No single dimension suffices for all purposes, because
the landscapes are genuinely different, as the separation results of
\Cref{ch:separations} will make precise.

The next chapter turns from combinatorial dimensions to a different family of
complexity measures: those based on \emph{compression} and \emph{sample
complexity}.

\section{The Dimensions of an Attention Class}
\label{sec:attention-dimensions}

The combinatorial dimensions of this chapter take concrete values on the attention classes of a transformer, and those values control its sample complexity, its online learnability, and its learning rate.  Fix a finite-head attention class $\mathcal{A}_{d,n_K}$ on $\R^d$ with $n_K$ keys.

\begin{prop}[The VC dimension of attention is finite]
\label{prop:attention-vc-dim}
A fixed-architecture soft-attention class has finite VC dimension, so the Sauer--Shelah bound (\cref{lem:sauer-shelah}) applies: its growth function is polynomial.\formal{FLT_Proofs.Complexity.VCDimension.VCDim}
\end{prop}

\begin{proof}
The class is a well-behaved VC-measurable target (\cref{prop:attention-pac}), and its concepts are cut out by a bounded number of smooth inequalities in finitely many parameters; such parametric families shatter no arbitrarily large set, so $\VC(\mathcal{A}_{d,n_K}) < \infty$.  Sauer--Shelah then bounds $\growth_{\mathcal{A}}(n) \leq (en/\VC)^{\VC}$.
\end{proof}

The other dimensions of this chapter then read off the rest of the attention class's profile.  Its \emph{Littlestone} dimension (finite or infinite) decides whether the class is online learnable and whether its universal rate is exponential or linear (\Cref{ch:universal}); finiteness is open (\cref{op:attention-dim-profile}).  Its \emph{star number} governs how cheaply an active learner can extract the head's behaviour by querying, the methodology of \Cref{ch:exact}.  And when attention routes over a label set of size $n_K \geq 3$ (a genuine multiclass head), its \emph{DS} dimension, not its Natarajan dimension, is what decides learnability, by the same thirty-year correction told above.  The single architecture is thus a point in every one of this chapter's dimension spaces at once, and the spaces disagree about it in exactly the ways the separation results of \Cref{ch:separations} predict.

\section{Open Problems}
\label{sec:ch10-open}

Each problem is anchored to a result above and tagged: a \emph{known unknown} (\textsc{ku}) is an articulated open question; an \emph{unknown known} (\textsc{uk}) is a result in the literature or the verified development not yet absorbed into a clean statement.  Verified handles are named where they exist.

\begin{openprob}[The dimension profile of attention, \textsc{uk}]
\label{op:attention-dim-profile}
\Cref{prop:attention-vc-dim} pins the VC dimension of a finite-head attention class as finite.  Compute the full dimension profile (VC, Littlestone, star number, and, for multiclass heads, DS) as functions of the architecture and score temperature, settling the regime questions of \Cref{ch:online,ch:universal} for attention in one stroke.
\end{openprob}

\begin{openprob}[The optimal compression of VC classes, \textsc{ku}]
\label{op:od-compression-dim}
The Sauer--Shelah lemma (\cref{lem:sauer-shelah}) converts finite VC dimension into a polynomial growth function.  Determine whether every class of VC dimension $d$ admits a sample compression scheme of size $O(d)$ (Littlestone--Warmuth), the open conjecture isolated in \leanname{FLT_Proofs.Complexity.Structures.Open_NoInfoCompressionStrengthening} and recurring from \Cref{ch:objects,ch:compression}.
\end{openprob}

\begin{openprob}[A combinatorial proof that DS characterizes multiclass learning, \textsc{uk}]
\label{op:ds-formal}
The DS characterization (\cref{thm:ds-char}) is cited, not verified.  The development types the Natarajan dimension and an \emph{all-labelings} shattering dimension (\leanname{StrongShatteringDim}, whose value is at most $\Ndim$) but \emph{not} the genuine Daniely--Shalev-Shwartz DS dimension, whose pseudo-cube definition gives the reverse bound $\Ndim \le \DSdim$.  Formalize the true DS dimension and the Brukhim et al.\ proof that $\DSdim(\mathcal{H}) < \infty$ iff $\mathcal{H}$ is multiclass PAC learnable.
\end{openprob}

\begin{openprob}[Closing the Natarajan--DS gap quantitatively, \textsc{ku}]
\label{op:natarajan-ds-gap}
\Cref{thm:natarajan-failure} shows $\Ndim$ can be $1$ while the class is unlearnable ($\DSdim = \infty$).  Characterize exactly how large the DS dimension can be as a function of the Natarajan dimension and the label-set size $|Y|$, and identify the structural feature of the label set that makes the gap appear.
\end{openprob}

\begin{openprob}[Fat-shattering rates for attention, \textsc{uk}]
\label{op:fat-attention}
Soft attention outputs real-valued scores before thresholding.  Compute the fat-shattering dimension $\fatshat_\gamma$ (\cref{def:fat-shattering}) of the real-valued attention-score class as a function of the margin $\gamma$ and the temperature, giving margin-based generalization bounds for attention via \cref{thm:fat-char}.
\end{openprob}

\begin{openprob}[A unifying dimension functor, \textsc{ku}]
\label{op:dimension-functor}
The chapter catalogs a dozen dimensions, each governing one paradigm.  Determine whether they are values of a single functor from concept classes (with label-preserving maps) to ordinals, recovering VC, Littlestone, DS, and fat-shattering by specializing a data-model and a scale parameter, or prove the dimensions are provably non-unifiable.
\end{openprob}

\begin{openprob}[The exact star-number--VC gap, \textsc{ku}]
\label{op:star-vc-gap}
\Cref{rmk:star-number} corrects the folklore: $\mathfrak{s} \geq \VC$ with an infinite gap for singletons.  Characterize the classes for which $\mathfrak{s}(\mathcal{H}) = O(\VC(\mathcal{H}))$, and determine the active-learning label-complexity speed-up that finiteness of the star number buys over the worst case, on top of \leanname{FLT_Proofs.Complexity.VCDimension.StarNumber}.
\end{openprob}

\begin{openprob}[Pseudodimension versus fat-shattering for neural classes, \textsc{uk}]
\label{op:pdim-fat}
\Cref{thm:fat-char} shows the pseudodimension alone does not characterize agnostic real-valued learnability; the scale-sensitive fat-shattering dimension does.  Determine, for the score class of a transformer, how far the pseudodimension over-estimates the learnable complexity, and whether the gap shrinks with depth.
\end{openprob}

\begin{openprob}[Ordinal dimensions and the slow regime, \textsc{uk}]
\label{op:ordinal-dim-slow}
\Cref{tab:other-dims} lists the ordinal VC and Littlestone dimensions (\leanname{FLT_Proofs.Complexity.Ordinal.OrdinalLittlestoneDim}).  Determine whether a single ordinal-valued dimension characterizes the boundary between the linear and arbitrarily-slow universal rates of \Cref{ch:universal}, completing the dimension story of the trichotomy.
\end{openprob}

\begin{openprob}[The dual VC dimension and attention routing, \textsc{uk}]
\label{op:dual-vc-attention}
The development types the dual class and proves a dual-to-primal shattering transfer (\leanname{FLT_Proofs.Complexity.DualVC.dual_shatters_imp_original_shatters}).  Interpret the dual VC dimension of an attention class (shattering by the \emph{keys} rather than the inputs) and determine what it measures about a transformer's routing capacity.
\end{openprob}

\subsection*{Counterfactual exercises: generalizing Figure~\ref{fig:shatter-halfplane}}

Each exercise perturbs the shattering picture.  Argue, in the figure's grammar (a $2^d$-cell grid in which each cell carries one labeling of a fixed $d$-point set together with the separating halfplane that realizes it, so that filling every cell with a witness exhibits the restriction $\mathcal{H}|_S$ as the full cube $\{0,1\}^d$, drawn solid where the kernel certifies the concept (the definitions \leanname{Shatters} and \leanname{VCDim} and the deep role \leanname{vc_characterization}) and dashed where the \emph{value} $\VC=3$, its generalization $\VC=d+1$, and the four-point obstruction are classical, cited not proved), that the perturbed picture subsumes the concrete one as an \emph{instance} (a class whose largest fillable grid is exhibited), a \emph{boundary} (one labeling or one point toggled so a cell becomes unrealizable), or a \emph{frontier} (a dashed, open or literature-only edge).

\begin{enumerate}[leftmargin=*,itemsep=3pt]
\item \textbf{The grid is the definition of shattering, not a fact about the plane.} Read each filled cell not as a colored triangle but as one element of the restriction. Show this is the figure's defining instance of the verified concept: a witness in every cell says exactly $\mathcal{H}|_S = \{0,1\}^{|S|}$, which is the definition of shattering (\cref{def:shattering}, \leanname{Shatters}), so the eight halfplanes are not decoration but the constructive content of ``$S$ is shattered''; the picture certifies \leanname{Shatters} cell by cell rather than asserting it.

\item \textbf{The value three is the largest fillable grid, and it is literature.} Read the whole figure as a single number. Show this is the instance that names the second object: $\VC(\mathcal{H})$ is the largest $|S|$ for which the grid fills (\cref{def:vc-dim}, \leanname{VCDim}), so the figure exhibits $\VC(\mathcal{H}) \ge 3$ by construction, but the matching upper bound $\VC = 3$ is the classical value (\cite{VapnikChervonenkis1971}, \cite{BlumerEhrenfeuchtHausslerWarmuth1989}), drawn dashed because the kernel types \leanname{VCDim} without proving the halfplane value; the concept is verified, the number is cited.

\item \textbf{The fourth point empties a cell.} Add a fourth point inside the convex hull of the other three and ask for the labeling that flips it against the rest. Show this is the figure's principal boundary: that single cell of the $2^4$ grid can never carry a separating line, because one point in the hull of the others forbids the opposite-label halfplane (Radon, \Cref{ex:halfplanes-shatter}), so the grid for $d=4$ has a permanently empty cell and $\VC = 3$; the obstruction is convex-geometric and classical, with no kernel theorem, so it is drawn dashed.

\item \textbf{Why filling the grid is what makes the class learnable.} Ask why the largest-fillable-grid number governs anything. Show this is the frontier that gives the figure its weight: the very finiteness the grid measures is what the Fundamental Theorem turns into PAC learnability, since $\mathcal{H}$ is PAC learnable if and only if $\VC(\mathcal{H}) < \infty$ (\cref{thm:fundamental}, \leanname{vc_characterization}, \leanname{fundamental_theorem}); the halfplane class, with its grid stalling at $d=4$, is therefore learnable, and this if-and-only-if is the one solidly verified edge attached to the picture.

\item \textbf{Shrink to intervals on the line.} Replace halfplanes in $\R^2$ by halflines, or by intervals, on $\R$. Show this is the instance one dimension down: a single threshold shatters one point but not two of the right pattern, so the largest fillable grid is the $2^1$ grid and $\VC = 1$ for halflines, while intervals fill the $2^2$ grid and $\VC = 2$; both values instantiate the same \leanname{VCDim} definition and are the classical geometric values (\cite{BlumerEhrenfeuchtHausslerWarmuth1989}), so the picture's $d$ is a free parameter, not fixed at three.

\item \textbf{Grow to axis-aligned rectangles.} Replace halfplanes by axis-aligned rectangles in $\R^2$. Show this is the instance that raises the number without changing the grammar: four points placed as the extreme north, south, east, and west witnesses fill the $2^4$ grid, while no five points do, so $\VC = 4$ for rectangles (\cite{BlumerEhrenfeuchtHausslerWarmuth1989}); the figure's $2^d$ grid is unchanged, only the witness shape (a rectangle in each cell instead of a line) and the largest fillable $d$ differ, both still instances of \leanname{VCDim}.

\item \textbf{Climb to halfspaces in $d$ dimensions.} Replace the plane by $\R^d$ and halfplanes by halfspaces. Show this is the instance that states the figure's general value: $d+1$ points in general position fill the $2^{d+1}$ grid and no $d+2$ do, so $\VC = d+1$ for halfspaces (\cite{BlumerEhrenfeuchtHausslerWarmuth1989}), with the $d=2$ picture as the case $\VC = 3$; the value $d+1$ is classical and not in the kernel (there is no halfspace dimension theorem), so it is dashed, and formalizing it together with the Radon upper bound is open.

\item \textbf{The grid never reaches the Littlestone grid.} Ask whether filling the static $2^d$ grid is the same as winning the adaptive online game. Show this is the boundary between two dimensions: shattering a set is non-adaptive (one fixed $S$, all labelings at once), so $\VC(\mathcal{H}) \le \Ldim(\mathcal{H})$ always (\cref{prop:vc-leq-ldim-dim}, \leanname{BranchWiseLittlestoneDim_ge_VCDim}), and the gap is infinite for thresholds, whose static grid fills only to $d=1$ yet whose adaptive mistake tree is unbounded; the figure draws the static witness, and the adaptive refinement is a strictly larger object certified by the same kernel.

\item \textbf{Swap the points and the halfplanes.} Read the figure with roles reversed: let the eight labelings shatter the family rather than the three points shatter the cube. Show this is the frontier of the dual view: the one-inclusion and dual-shattering reading swaps inputs for hypotheses, and the kernel types the dual class with a dual-to-primal shattering transfer (\cref{op:dual-vc-attention}, \leanname{FLT_Proofs.Complexity.DualVC.dual_shatters_imp_original_shatters}), but the self-dual value $\VC = d+1$ under this swap is not drawn and not formalized, so the dual reading of the grid is a dashed, literature-and-open edge.

\item \textbf{Is the grid the rank of a functor?} Read ``the largest finite set on which $\mathcal{H}$ realizes every labeling'' as a universal property. Show this is the figure's deepest frontier: $\VC$ would be the rank of a representable functor on finite sets, the largest free object (full cube) the class can hit, and the chapter's other dimensions (Littlestone, Natarajan \leanname{NatarajanDim}, fat-shattering) would be the same functor specialized by a data-model and a scale; whether a single such dimension functor exists, recovering all of them, or whether the dimensions are provably non-unifiable, is open (\cref{op:dimension-functor}), so this reading of the grid is drawn dashed.
\end{enumerate}

\subsection*{Counterfactual exercises: generalizing Figure~\ref{fig:growth-function}}

Each exercise perturbs the growth-function picture.  Argue, in the figure's grammar (a single curve $\growth_\mathcal{H}(n)$ that rides the $2^n$ reference ray exactly while the class still shatters, $n \le d = \VC(\mathcal{H})$, then peels permanently into the degree-$d$ polynomial $\sum_{i=0}^{d}\binom{n}{i}$ once $n$ passes $d$, drawn solid where the kernel certifies the object the curve traces (the definition \leanname{GrowthFunction} and the bound \leanname{sauer_shelah_quantitative}) and dashed where the asymptotic form $(en/d)^d$ and the immediate-and-permanent reading are classical, cited not proved), that the perturbed picture subsumes the concrete $d=3$ one as an \emph{instance} (a class whose peeling curve bends at its own $\VC$), a \emph{boundary} (one regime, point, or label toggled so the transition moves or vanishes), or a \emph{frontier} (a dashed, open or literature-only edge).

\begin{enumerate}[leftmargin=*,itemsep=3pt]
\item \textbf{The peeling curve is the growth function.} Read the solid curve not as a bound drawn around the class but as the quantity itself. Show this is the figure's defining instance of the verified object: the curve plots $\growth_\mathcal{H}(n) = \max_{|S|=n}|\mathcal{H}|_S|$ for a maximum class of VC dimension three (\cref{def:growth-function}, \leanname{GrowthFunction}), so its coincidence with $2^n$ up to $n=d$ and its descent afterward are the behavior of a single defined function, and the figure certifies \leanname{GrowthFunction} value by value rather than asserting two disconnected ceilings.

\item \textbf{The bend is the bound, and the bound is verified.} Read the moment the curve leaves the $2^n$ ray. Show this is the instance that names the second object: for $n > d$ the curve follows the Sauer--Shelah bound $\sum_{i=0}^{d}\binom{n}{i}$ (\cref{lem:sauer-shelah}, \leanname{sauer_shelah_quantitative}), a genuine kernel theorem and not the legacy existential bound, so this peeling edge of the figure is solidly verified, while the closed form $(en/d)^d$ named beside it is the classical simplification (\cite{Sauer1972}, \cite{Shelah1972}), drawn as annotation because no kernel declaration proves that form.

\item \textbf{Thresholds bend at $d=1$: a line.} Replace the $d=3$ class by halflines on $\R$. Show this is the instance one degree down: a single threshold shatters one point but not two, so $\VC = 1$, the curve rides $2^n$ only at $n \le 1$ and then follows the degree-one polynomial $\sum_{i=0}^{1}\binom{n}{i} = 1+n$, a straight line, peeling from $2^n$ at $n=2$; the same \leanname{GrowthFunction} and \leanname{sauer_shelah_quantitative} apply with $d=1$, so the figure's bend-degree is a free parameter, not fixed at three (\cite{BlumerEhrenfeuchtHausslerWarmuth1989}).

\item \textbf{Rectangles bend at $d=4$.} Replace the class by axis-aligned rectangles in $\R^2$. Show this is the instance that raises the degree without changing the grammar: four extreme witnesses fill the $2^4$ grid and no five do, so $\VC = 4$, the curve coincides with $2^n$ through $n=4$ and peels into a quartic $\sum_{i=0}^{4}\binom{n}{i}$ for $n>4$; only the transition point ($n=d+1=5$) and the terminal slope (degree four) move, both still instances of \leanname{sauer_shelah_quantitative} (\cite{BlumerEhrenfeuchtHausslerWarmuth1989}).

\item \textbf{A finite class flattens to a constant.} Replace the class by a finite $\mathcal{H}$ with $|\mathcal{H}| = N$. Show this is the degenerate instance: $\growth_\mathcal{H}(n) \le \min(2^n, N)$, so the curve rides $2^n$ until $2^n$ exceeds $N$ and is then pinned flat at $N$, the polynomial of degree $\VC \le \log_2 N$ collapsing to a constant ceiling; the same \leanname{GrowthFunction} bounds it, and the bend is here a hard horizontal cap rather than a polynomial arc, the smallest nontrivial case of the dichotomy.

\item \textbf{Infinite VC: the ray that never bends.} Take a class that shatters arbitrarily large sets ($\VC = \infty$), for instance all finite and cofinite subsets of $\N$. Show this is the principal boundary: $\growth_\mathcal{H}(n) = 2^n$ for every $n$, so the curve is the straight $2^n$ ray with no peel, the Sauer--Shelah hypothesis $\VC = d$ never holds for any finite $d$, and \cref{lem:sauer-shelah} simply does not apply; this is the case the Fundamental Theorem excludes, since $\mathcal{H}$ is not PAC learnable when $\VC = \infty$ (\cref{thm:fundamental}, \leanname{vc_characterization}).

\item \textbf{The transition sits exactly at $n=d+1$.} Ask where the bound first constrains the curve below $2^n$. Show this is the boundary that locates the phase change: at $n=d$ a set can still be shattered so $\growth$ may equal $2^d$, but at $n=d+1$ the side condition $d \le m$ of \leanname{sauer_shelah_quantitative} fires and $\sum_{i=0}^{d}\binom{n}{i} < 2^n$ for the first time, forcing the peel; the drop is immediate, with no $n$ of intermediate growth, and permanent, which is the classical content (\cite{Sauer1972}, \cite{Shelah1972}) the marker makes visible.

\item \textbf{$\growth = 2^n$ is the definition of shattering, read on the curve.} Ask why the curve equals $2^n$ exactly on $n \le d$ and not by approximation. Show this is the definitional edge: $\growth_\mathcal{H}(n) = 2^n$ holds iff $\mathcal{H}$ shatters some size-$n$ set (\cref{def:shattering}, \leanname{Shatters}), so the on-the-ray portion of the curve is not a bound being nearly tight but the literal statement that $\mathcal{H}|_S = \{0,1\}^n$ for some $S$; the coincidence with $2^n$ is a definition unfolding, not a theorem, and the figure shows shattering and its failure as one continuous picture.

\item \textbf{Is the bound tight from below?} Ask whether every class of VC dimension three grows as fast as the solid curve, or can stay strictly below it. Show this is the figure's tightness frontier: the kernel certifies only the upper bound (\leanname{sauer_shelah_quantitative}) and the solid curve plots a maximum class that attains it, but no declaration proves a matching lower bound, so whether a given class's growth reaches $\Theta(n^d)$ or drops further is classical for the extremal witness classes (\cite{Sauer1972}) and unformalized in general; this tightness is a dashed, open edge, and the degree it pins is what an optimal compression scheme of size $O(d)$ would also measure (\cref{op:od-compression-dim}).

\item \textbf{Is the bend-degree the rank of a graded object?} Read ``the polynomial that bounds $\growth$ has degree $\VC$'' as a statement about a graded structure. Show this is the figure's deepest frontier: the growth function would be the Hilbert or Poincar\'e series of a graded object attached to $\mathcal{H}$ whose dimension is $\VC$, so ``polynomial of degree $d$'' becomes ``the representable functor has rank $d$'', recovering the chapter's other dimensions (Littlestone, Natarajan \leanname{NatarajanDim}, fat-shattering) as the same series specialized by a data-model and a scale; whether one such dimension functor exists, with the bend-slope as its rank, or the dimensions are provably non-unifiable, is open (\cref{op:dimension-functor}), so this reading of the curve is drawn dashed.
\end{enumerate}

\subsection*{Counterfactual exercises: generalizing Figure~\ref{fig:dimension-hierarchy}}

Each exercise perturbs the dimension ladder.  Argue, in the figure's grammar (a magnitude ladder of combinatorial dimensions $\mathrm{StrongShattering} \le \Ndim \le \DSdim$, with the VC dimension as the binary specialization of $\Ndim$, drawn solid where the kernel certifies the object (the definitions \leanname{VCDim}, \leanname{NatarajanDim}, \leanname{StrongShatteringDim}, the verified rung \leanname{strongShatteringDim_le_NatarajanDim}, and the binary characterization \leanname{vc_characterization}) and dashed where the genuine DS dimension, the upper rung $\Ndim \le \DSdim$, the multiclass characterization, and the Natarajan failure are literature, cited not proved), that the perturbed picture subsumes the concrete ladder as an \emph{instance} (a label-set regime or class on which a rung collapses or is attained), a \emph{boundary} (one label-set size or one inequality toggled so a rung reverses or a characterization breaks), or a \emph{frontier} (a dashed, open or literature-only edge).

\begin{enumerate}[leftmargin=*,itemsep=3pt]
\item \textbf{Two labels collapse the ladder to the Fundamental Theorem.} Set $|Y| = 2$ and read the whole ladder. Show this is the figure's defining instance of the binary regime: when the label set has two elements the Natarajan dimension equals the VC dimension exactly (\cref{def:natarajan}, \leanname{NatarajanDim}), the strong-shattering dimension equals both, and the three rungs degenerate to the single binary characterization $\VC(\mathcal{H}) < \infty \Leftrightarrow \mathcal{H}$ is PAC learnable (\cref{thm:fundamental}, \leanname{vc_characterization}, \leanname{fundamental_theorem}); the ladder is the multiclass refinement of one verified equivalence, and the $|Y|=2$ slice is the only place the entire figure is solid.

\item \textbf{The lower rung is the figure's only solid multiclass edge.} Read the bottom inequality alone. Show this is the instance that names the verified handle: from any set on which every labeling is realizable one builds a two-colouring Natarajan witness, so $\mathrm{StrongShatteringDim}(\mathcal{H}) \le \Ndim(\mathcal{H})$ for $|Y| \ge 2$ (\leanname{strongShatteringDim_le_NatarajanDim}), a genuine kernel theorem proved by exhibiting the witness, not asserted; this is the single machine-checked rung of the multiclass ladder, and the figure draws it solid while everything above it is cited.

\item \textbf{The Natarajan dimension is the obvious lift, and it equals VC at the base.} Read the $\Ndim$ node as the natural generalization. Show this is the instance that states the figure's central object: the Natarajan dimension replaces ``all $2^d$ binary labelings'' by ``all $2^d$ labelings drawn from two opposing functions'' (\cref{def:natarajan}, \leanname{NatarajanDim}), which is exactly the VC dimension when $|Y|=2$ and a strict generalization otherwise; the figure draws it solid as a verified definition, leaving open only whether finiteness of it suffices for learning, which the next exercise denies.

\item \textbf{Finite Natarajan dimension does not imply learnability.} Toggle the question from ``is $\Ndim$ defined'' to ``does $\Ndim < \infty$ characterize learning''. Show this is the figure's principal boundary, the red dashed failure: there is a class with $\Ndim(\mathcal{H}) = 1$ that is not PAC learnable (\cref{thm:natarajan-failure}, \cite{DanielyShalevShwartz2014}), so the obvious dimension can be as small as possible and still miss unlearnability; the implication $\Ndim < \infty \Rightarrow$ learnable is drawn dashed red because it is false, the one edge the figure refutes rather than asserts.

\item \textbf{An infinite label set is where the obvious dimension fails.} Ask where the witness of the previous exercise lives. Show this is the boundary in the label-set sweep: the $\Ndim = 1$ unlearnable class requires $|Y| = \infty$ (\cref{thm:natarajan-failure}), because the Natarajan dimension reads only pairs of opposing functions and cannot see the combinatorial explosion of an unbounded label set; the gap is not informational but label-set-driven, and quantifying how large the DS dimension grows as a function of $|Y|$ when $\Ndim$ stays finite is open (\cref{op:natarajan-ds-gap}).

\item \textbf{The DS dimension is the correct characterization, and it is literature.} Replace the failed rung by the true one. Show this is the instance that names the third object: a class is PAC learnable if and only if its DS dimension is finite (\cref{thm:ds-char}, \cite{BrukhimCarmonDinurMoranYehudayoff2022}), the multiclass analogue of the Fundamental Theorem, with the DS dimension defined via pseudo-cubes in the one-inclusion hypergraph (\cref{def:ds-dim}); both the dimension and its equivalence are drawn dashed because the genuine pseudo-cube notion is not machine-checked, so the figure cites the characterization rather than verifying it.

\item \textbf{The kernel's handle is the reverse inequality, not the DS dimension.} Read the bottom rung against the top one. Show this is the boundary that separates the verified object from the literature object: the kernel proves $\mathrm{StrongShatteringDim} \le \Ndim$ (\leanname{strongShatteringDim_le_NatarajanDim}), whereas the genuine DS relationship is $\Ndim \le \DSdim$, the \emph{opposite} direction, and the DS dimension can be infinite while the Natarajan dimension is finite so no $\DSdim \le \Ndim$ bound can hold; the all-labelings strong-shattering dimension is therefore \emph{not} the DS dimension, and the figure must show the kernel's actual handle rather than silently substitute it for $\DSdim$.

\item \textbf{The full ladder has a verified floor and a literature ceiling.} Read all three rungs as one chain. Show this is the boundary between evidence grades: the complete ladder is $\mathrm{StrongShattering} \le \Ndim \le \DSdim$, of which only the left inequality is kernel-verified (\leanname{strongShatteringDim_le_NatarajanDim}) while the right inequality $\Ndim \le \DSdim$ is literature (\cite{BrukhimCarmonDinurMoranYehudayoff2022}); the line-style is the statement that multiclass learnability is currently half machine-checked, the binary characterization and the lower rung solid, the multiclass characterization and the upper rung dashed.

\item \textbf{Formalize the genuine DS dimension and its characterization.} Ask what it would take to make the dashed ceiling solid. Show this is the figure's headline frontier: formalizing the pseudo-cube DS dimension and the Brukhim et al.\ proof that $\DSdim(\mathcal{H}) < \infty$ iff $\mathcal{H}$ is multiclass PAC learnable would turn the dashed top rung and the dashed characterization into verified edges (\cref{op:ds-formal}), supplying the multiclass analogue of \leanname{vc_characterization}; until then the one-inclusion hypergraph and its pseudo-cubes (\cref{def:ds-dim}) are a typed-but-unproved target, and this is the chapter's principal open formalization.

\item \textbf{Is there a single dimension functor under the whole ladder?} Read the three dimensions as one object specialized three ways. Show this is the figure's deepest frontier: the VC, Natarajan, and DS dimensions would be values of a single functor from concept classes (with label-preserving maps) to ordinals, recovered by specializing a data-model and a label-set, with the one-inclusion hypergraph as the representable structure whose pseudo-cube rank is the learnability dimension; whether one such dimension functor exists, recovering all rungs of the ladder, or the dimensions are provably non-unifiable, is open (\cref{op:dimension-functor}), so this reading of the ladder is drawn dashed.
\end{enumerate}
